\begin{table}[ht]
\centering
\caption{Positioning of this paper relative to recent empirical work on the
post-ChatGPT shock to programming Q\&A platforms. A check mark indicates
that the cited study reports a substantive empirical result on the
corresponding margin. ``Task-level DDD'' refers to estimation via a
triple difference exploiting variation in task-level LLM substitutability;
``support-post funnel'' refers to decomposing the effect through
stages: questions $\to$ answered $\to$ accepted $\to$ accepted-not-closed
$\to$ accepted-non-negative-score $\to$ funnel-qualified post.}
\label{tab:literature_positioning}
\footnotesize
\setlength{\tabcolsep}{3.5pt}
\begin{adjustbox}{max width=\textwidth}
\begin{tabular}{lccccccc}
\toprule
Study & \makecell{Quantity\\decline} & \makecell{Quality /\\composition} &
      \makecell{Users /\\retention} & Answers &
      \makecell{Novelty /\\readability} &
      \makecell{Task-level\\DDD} &
      \makecell{Support-post\\funnel} \\
\midrule
del Rio-Chanona et al.\ (2024)   & \checkmark & \checkmark & & & & & \\
Burtch et al.\ (2024)            & \checkmark & \checkmark & \checkmark & & & & \\
Xue et al.\ (2026, ISR)          & \checkmark & \checkmark & \checkmark & \checkmark & & & \\
Shan \& Qiu (2025, ISR)          & \checkmark & \checkmark & & \checkmark & & & \\
\citet{quinngutt2024heterogeneous}  &            & \checkmark & & & \checkmark & & \\
\textbf{This paper}              & \checkmark & \checkmark & & & & \checkmark & \checkmark \\
\bottomrule
\end{tabular}
\end{adjustbox}
\end{table}
